% Limitations and Future Work. Built by the future pass from the instructions
% below. Do not process now. The prompt requires an honest account of what could
% not be run, had to be approximated, or was deferred in the codegen and
% execution steps, plus the simulation caveats and the legislative limitations.
% Do not shorten this; pair each limitation with a concrete remedy in a table.

[Write as connected prose plus one limitation-and-remedy table. Be candid and
specific; the credibility of the bill depends on stating plainly what the record
does and does not establish. Synthesise the two prior limitations sections and the
execution record; do not narrate them. Cover four groups.]

[Group 1, Codegen and execution limitations (what could not be run or was
approximated). From \texttt{VVUQ-01/final-paper/sections/limitations\_future.tex}
and \texttt{VVUQ-02/final-paper/sections/limitations\_future.tex}, and the
execution READMEs, state plainly: the agentic and on-prem LLM backends were not
live and ran on deterministic offline fallbacks (so intent generation used a
frozen reference plan); live web search and PDF ingestion were unreachable and
used offline stubs or size estimates; physics and robotics backends (for example
MuJoCo and Isaac) were not installed, so behaviors ran on standard-library
references rather than a full physics engine; local execution used a single Python
interpreter with cross-version coverage delegated to the CI matrix (3.10, 3.11,
3.12); the VVUQ-01 multibyte autochunk lost six characters on one oversized line
(reported, not repaired, in scope); and there was no independent gold-standard
validation reference in the first study, so one accept case was self-referential.
Note the two assurance gaps from VVUQ-01 (no independent reference; the rich input
corpus not yet wired into the executed pipeline). Cite \texttt{kawchak2026vvuq01}
and \texttt{kawchak2026vvuq02}.]

[Group 2, Simulation caveats. State that the Unitree H2-Surgical 1.0 is a clearly
labelled hypothetical 2030 platform, that every supporting number is a simulation
result, and that the four-entrant comparison is simulation against simulation;
note the time-dimension caveat on the human-surgeon baseline. State that a real
deployment would require IEC 80601-2-77, IEC 60601, ISO 13482, FDA SaMD Class III
clearance, IRB approval, and regulatory authorization, and that the bill does not
presume any of these. Cite \texttt{iec8060127}, \texttt{iec60601}, and
\texttt{iso13482}.]

[Group 3, This draft-paper step (process honesty). Note what this draft-paper step
itself could not do: pdflatex was not available in the container and the prompt
asked not to compile, so the LaTeX was validated structurally (balanced braces and
environments, every \texttt{\textbackslash cite} key present in the bibliography,
even dollar pairs, single-hyphen and section-symbol scans) and made compile-safe,
with the rendered PDF and any final page balancing left to Overleaf and the future
full-paper pass. No images are used. Cite \texttt{lamport1994latex}.]

[Group 4, Legislative limitations and future work. State that this is an
independent draft bill, not enacted law, and is not endorsed by CFR, ICH, FDA, or
any sponsor, CRO, site, IRB, regulator, or medical society; that it reflects one
author and needs the stakeholder input the policy memorandum names; and that the
state-law landscape is moving (the Texas committee substitute, the Florida 2026
bills), so citations should be refreshed at enactment. For future work, set the
forward claim that this verification-before-generation standard should extend
across the full breadth of oncology trials, with the remedies in the table
prioritised: supply an independent validation reference; wire the standards corpus
end to end; enable live backends; repair the multibyte chunker; and broaden the
simulation set toward real-world validation. Cite \texttt{kawchak2026bills},
\texttt{kawchak2026prediction}, and \texttt{fda2026realtime}.]

% Model limitation-and-remedy table (parent paths at top, abbreviated leaves):
% \begin{table}[h]\centering
% \begin{tabularx}{\textwidth}{L{5.0cm} Y}
% \toprule
% \textbf{Limitation (abbreviated source)} & \textbf{Concrete remedy} \\
% \midrule
% \multicolumn{2}{l}{\textit{Parent: papers/VVUQ-01/ and papers/VVUQ-02/}} \\
% LLM backends offline (execution READMEs) & Enable live backends; re-run with the same seed \\
% No independent reference (VVUQ-01 limits) & Supply a gold-standard reference; retest validation \\
% Single interpreter locally & Already covered by the CI 3.10--3.12 matrix \\
% Simulation against simulation (VVUQ-02 limits) & Stage toward physical validation under full standards \\
% \bottomrule
% \end{tabularx}
% \caption{Limitations from the codegen and execution steps, each with a remedy.}
% \end{table}
